\section{The filled scalar system and its a priori bounds}
\label{alg:section}

This section constructs the background on which the scalar deformation is
solved, proves its curvature identity, and obtains bounds for every smooth
solution in the continuation family.  The argument uses no energy condition
inside the filling.  In particular, different trapping signs at different
boundary components cause no incompatibility there.  Existence and higher
regularity of solutions are proved in the next section; none of those
conclusions is used in the estimates below.

Fix prepared exterior data as in the preceding section.  The properties used
here are the following: $K$ has compact support; in the distinguished end
$g_{ij}-\delta_{ij}=O_j(r^{2-n})$ for every fixed $j$;
$R_g=O(r^{-n-\beta})$ for a fixed $\beta>0$; the charged energy margin is
strict and bounded below by a positive multiple of $r^{-n-\beta}$; and
$H+\epsilon_a\tr_{S_a}K<0$ on each component $S_a$.  Set $k=n-1\geq3$.
All tensors, contractions, derivatives, and divergences in this section,
unless explicitly distinguished, use the background metric $g$.

\subsection{Filling and parameter conventions}

\begin{lemma}[A filling with long product collars]
\label{alg:filling}
For each sufficiently large integer $N$ the exterior has a smooth oriented
extension $(\mathcal M_N,g,K)$ without boundary, with the same single end,
having the following properties.
\begin{enumerate}[label=(\roman*)]
\item A fixed collar through every $S_a$, followed on its inner side by a
fixed transition to a product, has
$H+\epsilon_a\tr_{\mathrm{tan}}K<0$ on its cross-sections, with the normal
pointing toward the original exterior.
\item Behind that transition is a product cylinder of length
$L_N=C_0+C_1N$, with $K=-\epsilon_a L_0g$ there.  The cylinders terminate
at fixed collars of a compact filling; $K$ on the rest of that filling may
be arbitrary.
\item There are coordinate patches of a fixed size with uniform metric
ellipticity and uniform bounds for each fixed number of derivatives of
$g$ and $K$.  At most $C(1+N)$ such patches cover the added region and
the original compact core.  Its volume is at most $C(1+N)$.
\end{enumerate}
Here $L_0,C_0,C_1$ can be chosen in advance, with $L_0$ arbitrarily large.
The constants in (iii) are independent of $N$.
\end{lemma}

\begin{proof}
Cut the original exterior at a large coordinate sphere.  A second copy of
the resulting compact manifold, with its orientation reversed and its
spherical boundary capped by a ball, provides a compact filling of the
entire boundary $S$.  This construction does not require the individual
components of $S$ to bound separately.  Smoothly extend the original metric
and tensor a short distance behind $S$.  Strictness of each expansion
persists on a sufficiently short collar.  In its continuation, interpolate
the metric to a fixed product metric.  During that interpolation add
$-\epsilon_a Lg$ to $K$, with a nonnegative cutoff that becomes one before
the original strict collar is left.  Choose $L$ sufficiently large to
dominate the bounded mean curvatures of all intermediate leaves and the
remaining tangential trace.  One can then interpolate $K$ to
$-\epsilon_a L_0g$ while preserving the strict inequality.  This only
uses a fixed transition with bounded geometry.

Insert the product cylinder between this transition and the corresponding
boundary collar of the compact filling.  Prescribe the same pure trace on
the latter collar, and extend $K$ smoothly and arbitrarily farther into
the filling.  In particular, the prescriptions for opposite signs are
made on disjoint collars, and are not imposed on the common interior of
the filling.  There is no energy-condition requirement in that interior.
The fixed pieces have finite uniform atlases; translates of finitely many
product charts cover the cylinders.  This proves the last assertion and
the volume bound.  Increasing $C_0,C_1$ later to accommodate a prescribed
ramp does not change these local bounds.
\end{proof}

Extend $r$ to a smooth positive function, uniformly comparable to one on
the enlarged core and unchanged sufficiently far out.  This can be done
using the fixed collars and by making it constant on the added products.
Consequently every fixed negative power of $r$ is bounded above uniformly
in $N$, and bounded below on the enlarged core and any fixed additional
end annulus.

Fix $0<\varepsilon<1$ and a smooth nonincreasing function
$\vartheta:\R\to[0,1]$, equal to one on $(-\infty,0]$ and to zero on
$[1,\infty)$.  The order of choices is: prepared data and the fixed
geometric pieces, then $\varepsilon$, the exponents and small constants
below, then $N$, and finally the outer truncation radius.  Increasing
that radius will never change $N$ or any of the constants in its
equations.  We require
\begin{equation}
\begin{gathered}
 N\geq N_0>\max\{k,2/\varepsilon\},\qquad
 \ell=e^{-N\varepsilon},\qquad \eta=\ell^{3/2},\\
 p(t)=N\vartheta(Nt),\qquad
 l(t)=\exp\left(\int_0^t p(s)\,\dd s\right).
\end{gathered}
 \label{alg:profiles}
\end{equation}
Thus $0\leq p\leq N$, $l'=pl$, $l(t)=e^{Nt}$ for $t\leq0$, and
$l(t)\leq e$ for all $t$.  When $t\geq-\varepsilon$ one also has
$l\geq\ell$.  On $[-\varepsilon,-\varepsilon+1/N]$,
$\ell\leq l\leq e\ell$.

The notation $\mathcal P_N$ denotes a positive quantity bounded by
$C(1+N)^m$, where $C,m$ depend only on the fixed choices just described,
not on the solution or the truncation radius.  The polynomial may
increase from one occurrence to the next.  Thus an
$\exp(\mathcal P_N)$ bound means $\exp(C(1+N)^m)$, rather than a
constant with unrestricted dependence on $N$.  Constants for later
fixed-$N$ Schauder estimates need not have this restriction.

\subsection{Implicit variables and the curvature identity}

For smooth functions $f,Z$, define $t$ implicitly and then all other
quantities by
\begin{equation}
\begin{gathered}
 Z=t+\frac{1}{2k}\log D,\qquad D=1+l(t)^2|\nabla f|^2,
 \qquad w=lD^{-1/2}\nabla f,\qquad a_w=|w|,\qquad v=|w|^2,
 \\
 d=D^{-1}=1-v,\qquad \chi=\frac{kd}{k+pv},\qquad
 A_d=g^{-1}-w\otimes w,\qquad
 A_\chi=g^{-1}-\frac{k+p}{k+pv}w\otimes w,
 \\
 u=e^{(k-1)t}l\sqrt D,\qquad h=\eta f,\qquad
 H^f=lD^{-1/2}\nabla^2f,\qquad
 F=A_\chi^{ij}(K_{ij}+H^f_{ij}),
 \\
 V=uA_\chi\bigl(k\,\dd Z+K(w,\cdot)\bigr),\qquad
 \bar g=g+l^2\dd f\otimes\dd f,\qquad
 \widehat g=e^{2t}\bar g .
\end{gathered}
\label{alg:definitions}
\end{equation}
Here a positive contravariant tensor acts on covectors to give vectors,
and $|\xi|_A^2=A^{ij}\xi_i\xi_j$.  In a $g$-orthonormal frame with
last direction $e=\nabla f/|\nabla f|$, the eigenvalues of $A_d$ are
$(1,\ldots,1,d)$ and those of $A_\chi$ are
$(1,\ldots,1,\chi)$.  In particular,
\begin{equation}
 \frac{k}{k+N}A_d\leq A_\chi\leq A_d\leq g^{-1},
 \qquad 0<\chi\leq d\leq1.
 \label{alg:tensor-comparison}
\end{equation}

For fixed $\sigma=|\nabla f|$, the derivative with respect to $t$ of
the right side of the first equation in \eqref{alg:definitions} is
$1+pv/k\geq1$.  Its limits as $t\to\pm\infty$ are $\pm\infty$.
The equation therefore defines a unique smooth $t$ for every $(Z,\dd f)$,
including $\dd f=0$.  Differentiating spatially gives the useful identity
\begin{equation}
 k\,\dd Z=(k+pv)\,\dd t+H^f(w,\cdot).
 \label{alg:implicit-differential}
\end{equation}

For the next calculation, put $B=K+H^f$ and use the preceding frame to
write its transverse, mixed, and axial blocks as $T,M,j$.  Thus
$T=B|_{e^\perp}$, $M=B(e,\cdot)|_{e^\perp}$, $j=B(e,e)$, and
$\tr T=F-\chi j$.  Write $\dd t=y+x e^\flat$ and
$T^0=T-(\tr T)g|_{e^\perp}/k$, and set $s_p=p+(k-1)/2$.

\begin{proposition}[Exact curvature identity]
\label{alg:curvature-proposition}
Let $g$ be any smooth Riemannian metric and $K$ any smooth symmetric
tensor.  With $\tau=\tr_gK$,
$2\mu=R_g+\tau^2-|K|^2$, and $J=\Div(K-\tau g)$, the definitions
above give
\begin{equation}
 \frac12 e^{2t}R_{\widehat g}
 =\mu+J(w)+\mathcal T+w(F)-F\tau-u^{-1}\Div V,
 \label{alg:scalar-identity}
\end{equation}
where
\begin{equation}
\begin{split}
 \mathcal T={}&\frac12|T^0|^2+|M+s_pa_w y|^2
 -\frac{(k-1)^2}{4}v|y|^2
 +ks_p(|y|^2+dx^2)\\
 &-\frac{k-1}{2k}(F-\chi j)^2+\chi j^2-\chi Fj+F^2
 +2s_p\chi j a_wx .
\end{split}
\label{alg:remainder}
\end{equation}
The right side is smooth also at zero slope; the choice of $e$ there is
irrelevant.  No constraint inequality, flatness assumption, or restriction
on $p'$ is used in this identity.
\end{proposition}

\begin{proof}
We supply a tensorial derivation to retain all background-curvature terms.
For symmetric tensors set
$\mathsf B(A,B)=\langle A,B\rangle-(\tr A)(\tr B)$ and
$P_A=A-(\tr A)g$.  When such a tensor acts as an endomorphism, its first
index is raised with $g$.  Consider on $\R_z\times\mathcal M_N$ the stationary
Lorentz metric
\[
 \mathbf g=-l^2(\dd z-\dd f)^2+\bar g.
\]
The metric induced on $z=\mathrm{constant}$ is $g$.  Its lapse is
$U=l\sqrt D$ and its shift is $Uw=l^2\nabla f$.  With the future-normal
second-form convention of the theorem, its second form is
\begin{equation}
 b=-U^{-1}\sym\nabla(Uw)^\flat
   =-H^f-p(\dd t\otimes w^\flat+w^\flat\otimes\dd t).
 \label{alg:auxiliary-second-form}
\end{equation}
Let $\mathbf n=U^{-1}\partial_z-w$.  Contracted Gauss and the normal
Ricci equation give, respectively,
\[
 \mathbf R=R_g+(\tr b)^2-|b|^2-2\mathbf{Ric}(\mathbf n,\mathbf n),
 \qquad
 \mathbf{Ric}(\mathbf n,\mathbf n)
 =-\mathbf n(\tr b)-|b|^2+U^{-1}\Delta U.
\]
Because the coefficients are stationary,
$\mathbf n(\tr b)=-w(\tr b)$ and $\Div(Uw)=-U\tr b$.
Substitution proves
\begin{equation}
 \mathbf R=R_g+\mathsf B(b,b)
 -2U^{-1}\Div\bigl(\nabla U+(\tr b)Uw\bigr).
 \label{alg:stationary-scalar}
\end{equation}
Alternatively, the coordinate $z-f$ displays $\mathbf g$ as a static
warped product, so that
$\mathbf R=R_{\bar g}-2l^{-1}\Delta_{\bar g}l$.
Since $\dd V_{\bar g}=\sqrt D\,\dd V_g$ and $\bar g^{-1}=A_d$,
\[
 l^{-1}\Delta_{\bar g}l
  =U^{-1}\Div\bigl((U/l)A_d\dd l\bigr),\qquad
 \nabla U-(U/l)A_d\dd l=-U\bigl(b(w,\cdot)\bigr)^\sharp.
\]
The last equality follows by differentiating $U=l\sqrt D$ and using
\eqref{alg:auxiliary-second-form}; its components include all derivatives
of $l$.  Subtracting the two scalar formulas gives
\[
 \tfrac12R_{\bar g}=\tfrac12R_g+\tfrac12\mathsf B(b,b)
                      +U^{-1}\Div(UP_bw).
\]
Put $Q=K-b$.  The product rule and
$\sym\nabla(Uw)^\flat=-Ub$ imply
\[
 J(w)-U^{-1}\Div(UP_Kw)=\mathsf B(K,b).
\]
Expanding $\mathsf B(Q,Q)$ therefore yields
\begin{equation}
 \tfrac12R_{\bar g}
 =\mu+J(w)+\tfrac12\mathsf B(Q,Q)-U^{-1}\Div(UP_Qw).
 \label{alg:graph-scalar}
\end{equation}

The conformal scalar law, with dimension $n=k+1$, subtracts
$k\Delta_{\bar g}t+k(k-1)|\dd t|_{A_d}^2/2$ from
$R_{\bar g}/2$.  Moreover,
\[
 \Delta_{\bar g}t
 =U^{-1}\Div(UA_d\dd t)-p|\dd t|_{A_d}^2,
 \qquad u=e^{(k-1)t}U.
\]
Changing the divergence weight from $U$ to $u$ in
\eqref{alg:graph-scalar} consequently gives
\begin{equation}
\begin{split}
 \tfrac12e^{2t}R_{\widehat g}={}&\mu+J(w)
 +\tfrac12\mathsf B(Q,Q)+(k-1)P_Q(w,\nabla t)
 +ks_p|\dd t|_{A_d}^2\\
 &-u^{-1}\Div\bigl(u(P_Qw+kA_d\dd t)\bigr).
\end{split}
\label{alg:invariant-remainder}
\end{equation}
Equation~\eqref{alg:implicit-differential}, or its transverse and axial
components, gives
\begin{equation}
 u(P_Qw+kA_d\dd t)=V-uFw,
 \qquad
 u^{-1}\Div(uw)=F+(1-\chi)j-\tau+2s_pa_wx.
 \label{alg:flux-identities}
\end{equation}
For completeness, the transverse and axial components of
$P_Qw+kA_d\dd t$ are
$a_wM+(k+pv)y$ and $-a_w\tr T+kdx$, respectively; those of
$V/u-Fw$ are exactly the same.  The second equality follows by taking
the trace of \eqref{alg:auxiliary-second-form} and then differentiating
the weight $e^{(k-1)t}$.

It remains to expand the algebraic terms.  The blocks of $Q$ are
$T$, $M+pa_wy$, and $j+2pa_wx$.  Thus
\[
 \tfrac12\mathsf B(Q,Q)
 =\tfrac12|T^0|^2-\frac{k-1}{2k}(\tr T)^2
   +|M+pa_wy|^2-(\tr T)(j+2pa_wx),
\]
and
$P_Q(w,\nabla t)=a_w(M+pa_wy)\cdot y-a_w(\tr T)x$.
Insert these formulas and \eqref{alg:flux-identities} into
\eqref{alg:invariant-remainder}, and use $\tr T=F-\chi j$.
The mixed transverse terms complete to
$|M+s_pa_wy|^2-(k-1)^2v|y|^2/4$; the remaining terms are precisely
\eqref{alg:remainder}.  This proves \eqref{alg:scalar-identity}.
Formula~\eqref{alg:invariant-remainder} and
\eqref{alg:flux-identities} also express $\mathcal T$ by smooth
tensorial operations, proving the assertion at zero slope.
\end{proof}

\subsection{Coercivity with polynomial losses}

\begin{lemma}[Coercivity]
\label{alg:coercivity-lemma}
For $0\leq p\leq N$, $0\leq v<1$, the remainder is nonnegative and
\begin{equation}
 |T|^2+|M|^2+dj^2+F^2+|\dd t|_{A_d}^2
 \leq C_k(1+N)^2\mathcal T.
 \label{alg:coercivity}
\end{equation}
At $\dd t=0$, if
$E_0=|T^0|^2+|M|^2+\chi j^2+F^2$, then
\begin{equation}
 \tfrac12E_0\leq\mathcal T\leq
 \left(1+\frac{1}{2k}\right)E_0.
 \label{alg:stationary-coercivity}
\end{equation}
In addition, without any $N$ loss,
\begin{equation}
 \mathcal T\geq\frac{k^2-1}{4}|\dd t|_{A_d}^2.
 \label{alg:gradient-coercivity}
\end{equation}
\end{lemma}

\begin{proof}
Completing first the $x$ square and then the $F$ square gives the exact
formula
\begin{equation}
\begin{split}
 \mathcal T={}&\tfrac12|T^0|^2+|M+s_pa_wy|^2
 +\left(ks_p-\tfrac{(k-1)^2}{4}v\right)|y|^2\\
 &+ks_pd\left(x+\frac{\chi a_w}{kd}j\right)^2
 +\frac{k+1}{2k}\left(F-\frac{\chi}{k+1}j\right)^2
 +\frac{\chi\,[k(k+2)+v]}{2(k+1)(k+pv)}j^2 .
\end{split}
\label{alg:completed-squares}
\end{equation}
Indeed the last coefficient, before simplification, is
\[
 \chi-\frac{k\chi^2}{2(k+1)}
       -\frac{s_p\chi^2v}{kd}
 =\chi\frac{k-\frac{k^2}{2(k+1)}d-\frac{k-1}{2}v}{k+pv},
\]
and the numerator in the final fraction equals
$[k(k+2)+v]/[2(k+1)]$.

The coefficient of $|y|^2$ is at least $(k^2-1)/4+kp$.
The coefficient of $j^2$, divided by $d$, is at least
$c_k(1+N)^{-2}$.  In the variables $X=\sqrt d\,x$ and
$J_0=\sqrt d\,j$, the $x$ square is
$ks_p(X+a_wJ_0/(k+pv))^2$; its shear coefficient is bounded by $1/k$.
The shear from the $F$ square has coefficient
$\chi/((k+1)\sqrt d)\leq1/(k+1)$.  These observations control
$X^2,J_0^2,F^2$ with at most $C_k(1+N)^2$ loss.  Recovering $M$
costs at most $C_k(1+N)$, since
$s_p^2/(ks_p-(k-1)^2v/4)\leq C_k(1+N)$.
Finally $|T|^2=|T^0|^2+(F-\chi j)^2/k$.  This proves
\eqref{alg:coercivity}.

When $\dd t=0$, the $F,j$ part is
\[
 \frac{k+1}{2k}F^2-\frac{\chi}{k}Fj
     +\left(\chi-\frac{k-1}{2k}\chi^2\right)j^2.
\]
Put $z=\sqrt\chi\,j$ and bound the mixed term by
$(F^2+z^2)/(2k)$.  Since $0<\chi\leq1$, the claimed lower and
upper bounds follow, together with the separate $T^0,M$ terms.

For the sharper gradient bound, minimize over $M,F,j$, keeping $x,y$
fixed.  The transverse coefficient remaining is
$ks_p-(k-1)^2v/4\geq(k^2-1)/4$.  Put
\[
 B_v=\frac{k(k+2)+v}{2(k+1)},\qquad
 D_v=k-\frac{k^2}{2(k+1)}d.
\]
The axial coefficient divided by $d$ is
\[
 C(p,v)=\frac{ks_pB_v}{D_v+pv}.
\]
As $D_v-\frac{k-1}{2}v=B_v$, one has
$\partial_pC=kB_v^2/(D_v+pv)^2>0$.  At $p=0$ the inequality
$C(0,v)\geq(k^2-1)/4$ is equivalent to $D_v\geq kv$, which follows
from $D_v-kv=(k-k^2/(2(k+1)))d\geq0$.  This proves
\eqref{alg:gradient-coercivity}.
\end{proof}

\subsection{The equations, their homotopy, and the height estimate}

Choose once and for all
\begin{equation}
 \frac n4<b_1<\min\{n-2,n/2\},\qquad
 \rho_0=r^{-n-\beta},\qquad \rho_1=r^{-2b_1},\qquad
 \rho=c_*\rho_0>0.
 \label{alg:weights}
\end{equation}
The interval for $b_1$ is nonempty for $n\geq4$.  Choose $c_*$ so
small that $2\rho$ is below the prepared charged energy margin on the
original exterior.  In particular, $\rho$ has a positive lower bound
independent of $N$ on the enlarged core and fixed additional annuli.
The exponent restrictions give $b_1<k-1$, $2b_1<n$, and
$\rho_1^2\in L^1$ on the end.  All integrals of $\rho_0$,
$\rho_0^2$, and $\rho_1^2$ are $O(1+N)$.

Write $\mathcal M_{N,R}$ for the filled manifold truncated at $r=R$.
It has only this outer boundary.  The system to be solved is
\begin{equation}
\begin{gathered}
 F=h=\eta f,\qquad \Div V=u\Xi,\qquad f=Z=0\quad\text{on }r=R,\\
 \Xi=\delta_0\bigl(\mathcal T+\eta a_w|\nabla f|\bigr)+\rho
       -m_0C_N(\rho_0+v\rho_1),\qquad
 m_0=\vartheta\bigl(N(t+\varepsilon)\bigr).
\end{gathered}
\label{alg:system}
\end{equation}
The constant $\delta_0>0$ is fixed sufficiently small independently of
$N$; the floor lemma below specifies the choice.  The penalty $C_N$ is
then chosen polynomially large.  It is fixed before the truncation radius.

There are two parts of the continuation.  First replace $K$ by $sK$,
$0\leq s\leq1$, everywhere in $F,V,\mathcal T$, leaving the displayed
source otherwise unchanged.  At $s=0$ the trace equation implies $f=0$.
In the second part retain $K=f=0$ and replace $\Xi$ by $\lambda\Xi$,
$0\leq\lambda\leq1$.  At $\lambda=0$ the sole solution is $Z=0$.
We will prove uniform a priori estimates on the part of this family
where $\min t\geq-\varepsilon$; the floor lemma excludes equality.

\begin{lemma}[Height and outer boundary gradient]
\label{alg:trace-height}
Every smooth solution of the trace equation in the first homotopy part
satisfies, with constants independent of $N,R,s$,
\begin{equation}
 |h|\leq C_h,\qquad
 |h|\leq C\bigl(r^{-b_1}-R^{-b_1}\bigr)
 \quad(r_0\leq r\leq R),\qquad
 |\dd h|\leq CR^{-b_1-1}\quad(r=R),
 \label{alg:height-bounds}
\end{equation}
provided $R\geq2r_0$, where $r_0$ is fixed and $K=0$ on $r\geq r_0$.
For $R\geq R_{\min}(N)$ these estimates imply $t>-\varepsilon$ on
the outer sphere, without a prior assumption on the interior floor.
The same assertions hold in the second homotopy part, where $h=0$.
\end{lemma}

\begin{proof}
In terms of $h$ the trace equation reads
\begin{equation}
 \frac{A_\chi^{ij}\nabla_i\nabla_jh}
 {\sqrt{(\eta/l)^2+|\dd h|^2}}+\tr_{A_\chi}(sK)=h.
 \label{alg:height-equation}
\end{equation}
At a positive maximum, the first term is nonpositive and the second is
bounded by $n\sup|K|$; a negative minimum is treated with the opposite
sign.  The supremum of $|K|$ on the filling is uniform in $N$.
Taking $C_h>n\sup|K|$ proves the first assertion.  The same argument
with $K=0$ proves $h=0$ on the last homotopy part.

Let $B_R=C(r^{-b_1}-R^{-b_1})$.  At a contact with $h-B_R$, the
axis of $A_\chi$ is the radial-gradient direction, so that
\[
 \tr_{A_\chi}\nabla^2B_R
 =Cb_1r^{-b_1-2}
   \bigl((b_1+1)\chi-k+O(r^{2-n})\bigr)<0.
\]
The error is uniform for $0<\chi\leq1$; choose $r_0$ large enough
using $b_1<k-1$.  Choose $C$ so that $B_R\geq C_h$ on $r=r_0$
for every $R\geq2r_0$.  At a positive interior maximum of $h-B_R$,
the Hessian term in \eqref{alg:height-equation} would be negative
while $h>B_R\geq0$, a contradiction.  The lower barrier is $-B_R$.
Their one-sided boundary derivatives give the last estimate in
\eqref{alg:height-bounds}; tangential derivatives there vanish.

On $r=R$, $Z=0$ and $l\leq e$.  Hence
\[
 t=-\frac{1}{2k}\log(1+l^2|\dd f|^2)
 \geq-\frac{1}{2k}\log\bigl(1+C\eta^{-2}R^{-2b_1-2}\bigr).
\]
For example, in addition to $R\geq2r_0$, require
$C\eta^{-2}R^{-2b_1-2}<e^{k\varepsilon}-1$.
This gives $t>-\varepsilon/2$ on the outer sphere and defines a
permissible $R_{\min}(N)$.
\end{proof}

\subsection{Exclusion of the floor}

Define, for a symmetric tensor $B$,
\[
 \mathcal S(B)=\tfrac12\bigl((\tr_{A_d}B)^2
                         -|B|_{A_d\otimes A_d}^2\bigr).
\]

\begin{lemma}[The stationary algebraic comparison]
\label{alg:floor-algebra}
There are constants $c_1>0$ and $C_k>0$, independent of $N$, such
that at any point where $\dd t=0$,
\begin{equation}
 \mathcal T-\mathcal S(H^f)
 \geq c_1\mathcal T
        -C_k(1+N)^2\bigl(vF^2+|K|^2\bigr).
 \label{alg:floor-quadratic}
\end{equation}
The same constants apply with $K$ replaced by $sK$, $0\leq s\leq1$.
\end{lemma}

\begin{proof}
First put $B=H^f+K$.  Its blocks give
\[
 \mathcal S(B)=\frac{k-1}{2k}(F-\chi j)^2
 -\tfrac12|T^0|^2+dj(F-\chi j)-d|M|^2.
\]
Subtracting this expression from \eqref{alg:remainder} at $\dd t=0$
yields
\begin{equation}
\begin{split}
 \mathcal T-\mathcal S(B)={}&|T^0|^2+(1+d)|M|^2+F^2/k\\
 &+\left(\frac{k-2}{k}\chi-d\right)Fj
 +\left(\chi(1+d)-\frac{k-1}{k}\chi^2\right)j^2.
\end{split}
\label{alg:floor-subtraction}
\end{equation}
The last coefficient is at least $\chi$, while the absolute value of
the mixed coefficient is at most $d$.  Fix a small dimension-dependent
$\delta_k>0$.  If $(1+p)v\geq\delta_k$, Young's inequality in
the variables $F,\sqrt\chi j$, together with
\[
 \frac{d^2}{\chi}=\frac{d(k+pv)}k\leq1+N/k,
 \qquad v^{-1}\leq(1+N)/\delta_k,
\]
gives
$\mathcal T-\mathcal S(B)\geq\tfrac12E_0
 -C_k(1+N)^2vF^2$.
If $(1+p)v<\delta_k$, both $d$ and $\chi$ differ from one by at
most $C_k\delta_k$.  The quadratic form in
\eqref{alg:floor-subtraction} then differs, by at most
$C_k\delta_k(|M|^2+F^2+j^2)$, from
\[
 |T^0|^2+2|M|^2+(F-j)^2/k+j^2.
\]
The least eigenvalue of the last $F,j$ form is at least $1/(k+2)$.
Choose $\delta_k$ small enough to preserve half this lower bound.
Combining the cases proves a fixed positive $E_0$ lower bound with
the stated error.

Finally, polarization of $\mathcal S$ and
$d^2/\chi\leq1+N/k$ give
\[
 |\mathcal S(B)-\mathcal S(H^f)|
 \leq C_k\sqrt{1+N}\sqrt{E_0}|K|+C_k|K|^2.
\]
Apply Young's inequality, using half the positive $E_0$ lower bound,
and then \eqref{alg:stationary-coercivity}.  This proves the result.
All steps are unchanged on replacing $K$ by $sK$.
\end{proof}

\begin{lemma}[Floor exclusion]
\label{alg:floor}
There is a choice $0<\delta_0<\min\{c_1,1\}/2$ and a polynomial
$C_N$ such that, for $R\geq R_{\min}(N)$, no smooth solution in
either part of the homotopy has $\min t=-\varepsilon$.
\end{lemma}

\begin{proof}
Lemma~\ref{alg:trace-height} excludes a boundary minimum at this value.
At an interior minimum $\dd t=0$, $\nabla^2t\geq0$.  View $\bar g$
as the metric of the graph $z=f(x)$ in the Riemannian warped product
$g+l^2\dd z^2$.  At this point its second form is $H^f$.  The ambient
scalar curvature is $R_g-2p\Delta t$.  With vertical unit vector
$e_z=l^{-1}\partial_z$ the graph unit normal is
$\sqrt d\,e_z-a_we$, and the ambient normal Ricci curvature equals
\[
 v\Ric_g(e,e)-vp\nabla^2t(e,e)-dp\Delta t.
\]
The Gauss and conformal formulas therefore give the exact equality
\begin{equation}
 \tfrac12 e^{2t}R_{\widehat g}
 =\tfrac12R_g-v\Ric_g(e,e)+\mathcal S(H^f)
 -vp\tr_{e^\perp}\nabla^2t-k\tr_{A_d}\nabla^2t.
 \label{alg:floor-gauss}
\end{equation}
The last two terms are nonpositive.  Compare this upper bound with
\eqref{alg:scalar-identity}, use $F=h$ and
$w(F)=\eta a_w|\nabla f|$, and apply
Lemma~\ref{alg:floor-algebra}.  It follows that
\begin{equation}
 u^{-1}\Div V
 \geq c_1\mathcal T+\eta a_w|\nabla f|
                  -P_N(\rho_0+v\rho_1)
 \label{alg:floor-divergence}
\end{equation}
for some polynomial $P_N$ independent of $R$ and the homotopy parameter.
Here are all the error estimates used in this conclusion.  The difference
$\mu-R_g/2=(\tau^2-|K|^2)/2$, $J(w)$, and $F\tau$ is uniformly
bounded on the enlarged core and vanishes on the end, so is bounded by
$C\rho_0$.  On the end $|\Ric_g|=O(r^{-n})\leq C\rho_1$ because
$2b_1<n$.  The height barrier gives $vF^2\leq Cv\rho_1$ there;
on the core its bound is absorbed by $C\rho_0$.  The same bounds hold
for $sK$, including its momentum density.

At $t=-\varepsilon$, $m_0=1$.  Choose $\delta_0$ as in the statement,
and then take, for instance, $C_N\geq P_N+c_*+1$.
Subtracting the right side of \eqref{alg:system} from
\eqref{alg:floor-divergence} gives a strictly positive quantity,
contradicting that equation.  This choice is polynomial in $N$.

In the last homotopy part, the height comparison gives $K=f=0$,
$t=Z$, $w=0$, and $\mathcal T=ks_p|\dd t|^2$.  At a floor minimum
the left side is $k\Delta t\geq0$, whereas for $\lambda>0$ the
right side is $\lambda(\rho-C_N\rho_0)<0$.  For $\lambda=0$ the
equation is $\Div(ku(t)\nabla t)=0$ with zero boundary data.
Testing by $t$ gives $\nabla t=0$, hence $t=0$.  This finishes the
homotopy check.
\end{proof}

\subsection{High levels and graph Sobolev estimates}

\begin{proposition}[Distortion bound]
\label{alg:distortion-proposition}
With the choices above, every smooth homotopy solution on
$\mathcal M_{N,R}$ with $\min t> -\varepsilon$ satisfies
\begin{equation}
 -\varepsilon<t\leq Z\leq\mathcal P_N,\qquad
 \ell\leq l\leq e,\qquad
 |f|+|\nabla f|\leq\exp(\mathcal P_N).
 \label{alg:distortion}
\end{equation}
The constants are independent of $R\geq R_{\min}(N)$ and of the
homotopy parameter.
\end{proposition}

\begin{proof}
We give the full high-level argument for the first homotopy part and write
$K$ for $sK$ throughout it.  Put $q_D=\sqrt D$ and
$a(t)=e^{(k-1)t}l$, so that
\begin{equation}
 u=a(t)q_D=l e^{kZ-t},\qquad
 \eta a_w|\nabla f|=\frac\eta l(q_D-q_D^{-1}),\qquad
 q_D=e^{k(Z-t)}.
 \label{alg:highlevel-factors}
\end{equation}
The penalty can be nonzero only on
$-\varepsilon<t<-\varepsilon+1/N$; its absolute size is at most
$\mathcal P_N$.  Equations~\eqref{alg:profiles} and
\eqref{alg:highlevel-factors} show that a positive threshold
$Z_0=\mathcal P_N$ can be chosen such that on $Z>Z_0$
\begin{equation}
 \Xi\geq\delta_0\mathcal T+\rho
                +\tfrac12\delta_0\eta a_w|\nabla f|.
 \label{alg:highlevel-sign}
\end{equation}
For example, on the penalty band it suffices to make
$(\eta/e)(e^{kZ}-1)$ larger than
$2C_N\sup(\rho_0+\rho_1)/\delta_0$; the logarithm of this requirement
is bounded by $C(1+N)+C\log(1+C_N)$.  Off the penalty band the
assertion is immediate.

Let $\zeta$ be smooth, nondecreasing, zero up to $Z_0$, and one
above $Z_0+1$, with bounded derivative.  Testing the divergence equation
by $\zeta(Z)$ is legitimate since $Z=0$ at the outer boundary.
Writing $A=A_\chi$ and $B_K=K(w,\cdot)$, we obtain
\[
 \int\zeta(Z)u\Xi\,\dd V_g
 =-\int\zeta'(Z)u
       \bigl(k|\dd Z|_A^2+\langle\dd Z,B_K\rangle_A\bigr)\,\dd V_g
 \leq C\int_{\{Z_0<Z<Z_0+1\}\cap\supp K}u|K|^2\,\dd V_g.
\]
On the last set $u\leq e\exp(k(Z_0+1)+\varepsilon)$.
The support has volume $O(1+N)$.  Using
\eqref{alg:highlevel-sign} proves
\begin{equation}
 \int_{\{Z>Z_0+1\}}u
    \bigl(\rho+\eta a_w|\nabla f|+\mathcal T\bigr)\,\dd V_g
 \leq\exp(\mathcal P_N).
 \label{alg:highlevel-integral}
\end{equation}

Let $\mathcal K_N$ be the enlarged core and a fixed additional annulus
containing $\supp K$ in its interior.  On this set $\rho\geq c>0$,
uniformly.  Since $l\geq\ell\geq\eta$ and $q_D\geq1$,
\[
 u\bigl(\rho+\eta a_w|\nabla f|\bigr)
 =q_De^{(k-1)t}\bigl(l\rho+\eta(q_D-q_D^{-1})\bigr)
 \geq c'\eta e^{(k-1)t}q_D^2.
\]
For any fixed $0<\alpha\leq k-1$, $t\geq-\varepsilon$ implies
$e^{\alpha Z}q_D=e^{\alpha t}q_D^{1+\alpha/k}
 \leq C e^{(k-1)t}q_D^2$.
On the complementary level set $Z\leq Z_0+1$,
$q_D\leq\exp(k(Z_0+1+\varepsilon))$.  Thus
\eqref{alg:highlevel-integral} and the core volume bound give, in
particular, the fixed-exponent seed
\begin{equation}
 \int_{\mathcal K_N}e^Z\sqrt D\,\dd V_g
 \leq\exp(\mathcal P_N).
 \label{alg:seed-integral}
\end{equation}
We chose $\alpha=1$; this exponent does not deteriorate as $N$ grows.

We next prove the local Sobolev inequality needed to turn
\eqref{alg:seed-integral} into a supremum estimate.  Use the graph of
$f$ in the ordinary product $(\mathcal M_N\times\R,g+\dd z^2)$.
Over a base patch its measure and inverse metric are
\[
 \dd V_\Gamma=S_f\,\dd V_g,\qquad
 S_f=\sqrt{1+|\nabla f|^2},\qquad
 A_\Gamma=g^{-1}-(1-S_f^{-2})e\otimes e.
\]
Both $\dd V_\Gamma$ and $A_\Gamma$ are comparable to
$\dd\mu_D=\sqrt D\,\dd V_g$ and $A_\chi$, respectively, with
$\exp(\mathcal P_N)$ factors.  In fact
\[
 \min\{1,l\}\leq q_D/S_f\leq\max\{1,l\},\qquad
 \frac{S_f^{-2}}\chi
 =\frac{1+l^2|\nabla f|^2}{1+|\nabla f|^2}\frac{k+pv}{k}.
\]
The second ratio lies between $\ell^2$ and $e^2(1+N/k)$.

The scalar mean curvature of this graph in the product is
\[
 H_\Gamma=\frac{q_D}{lS_f}
  \left(\tr(T-K|_{e^\perp})+S_f^{-2}(j-K(e,e))\right).
\]
The coefficients satisfy
\[
 \frac{q_D}{lS_f}\leq\max\{1,l^{-1}\},\qquad
 \frac{q_D}{lS_f}\frac{S_f^{-2}}{\sqrt d}
 =\frac{q_D^2}{lS_f^3}\leq\max\{l,l^{-1}\}.
\]
Consequently \eqref{alg:coercivity} gives
\begin{equation}
 |H_\Gamma|\leq\exp(\mathcal P_N)(1+\sqrt{\mathcal T}).
 \label{alg:graph-mean-curvature}
\end{equation}
Notice that this uses only $\sqrt d\,j$ and the transverse Hessian
block, not an unweighted bound for the whole Hessian.

For clarity concerning the ambient geometry in this step, each of the
finitely many fixed base models used in Lemma~\ref{alg:filling} can
be extended beyond the patch to a closed smooth manifold.  Choose a
smooth Euclidean isometric embedding of that fixed manifold
\cite{Nash1956}, and take its product with the line.  The bound on its
second fundamental form is fixed; it is the same on all translated
copies along a cylinder.  Its contribution to the Euclidean mean
curvature vector of the graph is at most $n$ times that bound, because
the horizontal projection of a graph unit vector has length at most
one.  Thus \eqref{alg:graph-mean-curvature} holds also for the Euclidean
mean curvature vector.  A compactly supported base test lifts to a
compactly supported graph test for each individual smooth $f$.

The Michael--Simon inequality \cite{MichaelSimon1973}, applied to
$|\varphi|^{2(n-1)/(n-2)}$ and followed by H\"older's inequality,
therefore gives
\begin{equation}
 \|\varphi\|_{L^{2n/(n-2)}(\dd\mu_D)}^2
 \leq\exp(\mathcal P_N)
 \int\bigl(|\dd\varphi|_{A_\chi}^2
                    +(1+\mathcal T)\varphi^2\bigr)\,\dd\mu_D
 \label{alg:graph-sobolev}
\end{equation}
for compactly supported tests in each patch.  The Euclidean theorem has
no slope or graph-area hypothesis.  The comparisons just established
check precisely the metric and mean-curvature terms used here.

It remains to supply a compatible energy estimate.  Globally
$\Xi\geq\delta_0\mathcal T-\mathcal P_N$.  Also
\[
 \dd\log a=(k-1+p)\dd t,\qquad
 |\dd\log a|_{A_\chi}\leq\mathcal P_N\sqrt{\mathcal T}
\]
by \eqref{alg:coercivity}.  Test the divergence equation by
$\psi^2e^{2sZ}/a$, where $\psi$ is a compactly supported base cutoff.
The weight $a$ cancels from $u=a\sqrt D$.  After moving the positive
principal term to the left, this gives
\begin{align*}
 &\int\psi^2e^{2sZ}
       (\delta_0\mathcal T+2ks|\dd Z|_A^2)\,\dd\mu_D\\
 &\quad\leq\mathcal P_N\int\psi^2e^{2sZ}\,\dd\mu_D
 +\int e^{2sZ}\psi^2\langle\dd\log a,k\dd Z+B_K\rangle_A\,\dd\mu_D\\
 &\qquad-2s\int e^{2sZ}\psi^2\langle\dd Z,B_K\rangle_A\,\dd\mu_D
 -2\int e^{2sZ}\psi\langle\dd\psi,k\dd Z+B_K\rangle_A\,\dd\mu_D.
\end{align*}
The logarithmic-weight terms cost at most
$\delta_0\psi^2\mathcal T/2+
 \mathcal P_N\psi^2(|\dd Z|_A^2+|K|^2)$ by Young's inequality.
The term multiplied by $s$ costs a fixed fraction of
$ks\psi^2|\dd Z|_A^2$ plus $Cs\psi^2|K|^2$.  The cutoff terms cost
another such fraction plus $C(1+s^{-1})|\dd\psi|_g^2+C\psi^2|K|^2$.
It follows, on choosing a polynomial $s_*(N)$ sufficiently large,
that for every $s\geq s_*(N)$,
\begin{equation}
 \int\psi^2e^{2sZ}
      (\mathcal T+s|\dd Z|_{A_\chi}^2)\,\dd\mu_D
 \leq\mathcal P_N(1+s)
       \int e^{2sZ}(\psi^2+|\dd\psi|_g^2)\,\dd\mu_D.
 \label{alg:weighted-energy}
\end{equation}
All absorption thresholds are polynomial in $N$.  No derivative of
$A_\chi$ is estimated in this test.

Set $W=e^Z$, $\kappa_*=n/(n-2)$, and $p_0=2s_*(N)>1$.
Insert $\varphi=\psi e^{sZ}$ into
\eqref{alg:graph-sobolev} and use \eqref{alg:weighted-energy}.
For two nested base patches whose gap is $b\in(0,1]$, the resulting
estimate, with $p=2s\geq p_0$, is
\begin{equation}
 \|W\|_{L^{p\kappa_*}(\mathrm{inner},\dd\mu_D)}
 \leq\bigl[\exp(\mathcal P_N)(1+p)^3b^{-2}\bigr]^{1/p}
          \|W\|_{L^p(\mathrm{outer},\dd\mu_D)}.
 \label{alg:moser-step}
\end{equation}
Iteration with $p_i=p_0\kappa_*^i$ and gaps proportional to
$b2^{-i}$ is legitimate: both $\sum_i1/p_i$ and $\sum_i i/p_i$
converge.  It gives
\begin{equation}
 \sup_{\mathrm{inner}}W
 \leq A_Nb^{-C_n}\|W\|_{L^{p_0}(\mathrm{outer},\dd\mu_D)},
 \qquad A_N\leq\exp(\mathcal P_N).
 \label{alg:moser-sup}
\end{equation}

We record explicitly why starting at the polynomially large $p_0$
does not spoil the bound obtained from \eqref{alg:seed-integral}.
Let $J_*\geq1$ bound that integral and take increasing patches with
successive gaps $b_j=b_*2^{-j-1}$, all inside a fixed larger patch.
If $m_j=\log\max\{1,\sup W\}$ on patch $j$, interpolation gives
\[
 \|W\|_{p_0}\leq(\sup W)^{1-1/p_0}J_*^{1/p_0},\qquad
 m_j\leq a+cj+(1-1/p_0)m_{j+1},
\]
where
$a=\log A_N+(\log J_*)/p_0+C_n\log(2/b_*)$ and
$c=C_n\log2$ can be increased to be positive.  The supremum on the
largest compact patch is finite for each individual smooth solution.
Iterating the recurrence and letting its length tend to infinity thus
gives
\[
 m_0\leq p_0a+c(p_0^2-p_0)\leq\mathcal P_N.
\]
This last estimate is polynomial even though $1-1/p_0$ tends to one.
The patches covering $\mathcal K_N$ have uniform sizes and geometry,
so the conclusion is uniform on that whole set.

Outside $\mathcal K_N$ one has $K=0$ on a neighborhood.  A larger
interior maximum of $Z$ above $Z_0$ would have
$\Div V=ku\tr_{A_\chi}\nabla^2Z\leq0$, contrary to
\eqref{alg:highlevel-sign}; the outer boundary has $Z=0$.
This proves the global upper bound for $Z$ in the first homotopy part.
In the last part $K=f=0$, $t=Z$, and the same high-maximum argument
applies whenever $\lambda>0$.  At $\lambda=0$ the solution is zero,
as already proved.

Finally $t\leq Z$ follows from the implicit relation,
$q_D=e^{k(Z-t)}\leq\exp(\mathcal P_N)$ follows from the floor,
and $|\nabla f|\leq q_D/l\leq q_D/\ell$.
Lemma~\ref{alg:trace-height} gives $|f|\leq C_h/\eta$.
These are the remaining assertions of \eqref{alg:distortion}.
\end{proof}

The estimates in this section are statements about every smooth solution,
uniformly in the outer truncation, and are therefore available before
continuation or exhaustion.  Their parameter order is essential:
$\delta_0$ is fixed first, $C_N$ is a polynomial penalty, and the
distortion bounds are exponential in a polynomial.  The subsequent
existence proof may use unrestricted constants at a fixed $N$, while
the later mass estimate will use only the polynomial losses proved here.
